\begin{figure*}[!ht]\centering
  \includegraphics[width=\textwidth]{Fig2}
  \caption{\textbf{Spatial deconvolution resolves seven niche-domain archetypes.}
(\textbf{a}) RCTD lineage composition of each domain archetype (56 sections).
(\textbf{b}) Spatial matrix: one representative section per stage (rows) against
H\&E, RCTD at the 6-lineage and 39-subtype level, the four spatial programmes
(dominant programme per spot) and the niche-domain archetype (columns). Scale bars are
calibrated per section from the nominal 100-$\mu$m spot spacing.
(\textbf{c}) RCTD subtype composition of each domain archetype.
(\textbf{d}) Domain composition per stage; D3 is confined to invasive disease
(17\% of spots, 0\% in all precursor stages) whereas the lymphoid domain D7 occurs in all
five stages.
(\textbf{e}) Relative make-up and absolute level of the four spatial programmes per stage.
\textbf{This panel is a negative control}: before depth matching the programme scores
essentially mirror sequencing depth, and after matching into depth deciles the stage
ordering is identical across all ten deciles but does not follow disease progression. The
fourth programme is the lipid-associated macrophage (LAM) signature \citep{cheng2021trem2lam},
which we report alongside three exhaustion axes; it is not an exhaustion signature in the
T-cell sense.
(\textbf{f}) Univariable Cox regression of each domain's marker signature on overall and
disease-free survival (TCGA-LUAD); p\,$<$\,0.05 in red. The D3 row uses the unadjusted
signature (see main text).
(\textbf{g}) Marker-gene dot plot for the seven domains.}\label{fig:2}
\end{figure*}
